% Shared innovation framework used across all brochures.
% Two-page spread: the first page teaches the three-pillar test; the second page
% translates the pillars into research, IP and industry-ready capability.
\SetAccent{149A98}{EAF7F6}
\CompactHero{The innovation framework}{Three pillars of innovation}{An idea becomes strategically interesting when it is \textbf{unique}, \textbf{hard to do} and \textbf{value adding} at the same time. The programs build student capability across all three.}

\begin{minipage}[t]{.318\textwidth}
\begin{SoftCard}
{\fontsize{8.2}{9.8}\selectfont\bfseries\color{Accent}01 \quad UNIQUENESS\par}\vspace{1.2mm}
{\fontsize{11.6}{13.4}\selectfont\bfseries\color{BroInk}Meaningfully different\par}\vspace{1.5mm}
\TightText{Students learn to ask whether the problem, insight or solution is genuinely different from what already exists - not merely a new implementation of a familiar idea.}
\vspace{2mm}\SmallCaps{Skills built}\par\vspace{.8mm}
\TightText{Problem discovery, reframing, prior-art awareness, assumption breaking, divergent thinking and TRIZ.}
\end{SoftCard}
\end{minipage}\hfill
\begin{minipage}[t]{.318\textwidth}
\begin{GoldCard}
{\fontsize{8.2}{9.8}\selectfont\bfseries\color{BroGold!80!black}02 \quad HARD TO DO\par}\vspace{1.2mm}
{\fontsize{11.6}{13.4}\selectfont\bfseries\color{BroInk}Technically defensible\par}\vspace{1.5mm}
\TightText{Innovation needs technical or execution depth. Students learn to work with constraints, contradictions, evidence and experiments so the idea is more than a superficial concept.}
\vspace{2mm}\SmallCaps{Skills built}\par\vspace{.8mm}
\TightText{Systems thinking, technical trade-offs, experimentation, prototyping, non-obviousness and execution discipline.}
\end{GoldCard}
\end{minipage}\hfill
\begin{minipage}[t]{.318\textwidth}
\begin{BroCard}
{\fontsize{8.2}{9.8}\selectfont\bfseries\color{Accent}03 \quad VALUE ADDING\par}\vspace{1.2mm}
{\fontsize{11.6}{13.4}\selectfont\bfseries\color{BroInk}Creates real value\par}\vspace{1.5mm}
\TightText{A clever idea is not automatically valuable. Students learn to connect innovation to a real user, institution, industry or societal need and define the change it should create.}
\vspace{2mm}\SmallCaps{Skills built}\par\vspace{.8mm}
\TightText{Stakeholder insight, outcome framing, evidence, adoption thinking, value articulation and storytelling.}
\end{BroCard}
\end{minipage}

\CardGap
\begin{BroCard}
\SectionTitle{Why all three have to come together}
\renewcommand{\arraystretch}{1.38}
\begin{tabularx}{\linewidth}{>{\RaggedRight\arraybackslash\bfseries\color{BroInk}\fontsize{8.8}{10.8}\selectfont}p{55mm} >{\RaggedRight\arraybackslash\color{BroMuted}\fontsize{8.8}{10.8}\selectfont}X}
Unique + valuable, but easy to copy & Interesting, but difficult to defend or sustain.\\
Unique + hard, but low value & Technically impressive, but unlikely to matter outside the project.\\
Hard + valuable, but not unique & Useful engineering, but limited differentiation or innovation claim.\\
\textcolor{Accent}{Unique + hard + value adding} & \textbf{A stronger foundation for defensible innovation and meaningful impact.}\\
\end{tabularx}
\end{BroCard}

\CardGap
\begin{BroCard}
\SectionTitle{How I help students build the three muscles}
\begin{minipage}[t]{.19\linewidth}\SmallCaps{Discover}\par\TightText{Find tensions, unmet needs and consequential problems.}\end{minipage}\hfill
\begin{minipage}[t]{.19\linewidth}\SmallCaps{Challenge}\par\TightText{Question assumptions and seek a genuinely different angle.}\end{minipage}\hfill
\begin{minipage}[t]{.19\linewidth}\SmallCaps{Build}\par\TightText{Work through constraints, contradictions and technical depth.}\end{minipage}\hfill
\begin{minipage}[t]{.19\linewidth}\SmallCaps{Validate}\par\TightText{Test evidence, stakeholder value, utility and feasibility.}\end{minipage}\hfill
\begin{minipage}[t]{.19\linewidth}\SmallCaps{Communicate}\par\TightText{Turn the work into a research, patent or industry story.}\end{minipage}
\end{BroCard}
\CardGap
\begin{DarkCard}{\color{BroGold}\fontsize{8.2}{9.8}\selectfont\bfseries THE CORE HABIT\par}\vspace{1mm}{\color{white}\fontsize{10.2}{12.3}\selectfont\bfseries Do not ask only, ``Can we build it?'' Ask: ``Is it different? Is it difficult enough to be defensible? Does it create enough value to matter?''}\end{DarkCard}
\ProgramFooter{UNIQUENESS \textbullet\ HARD TO DO \textbullet\ VALUE ADDING \textbullet\ IMPACT}
\newpage

\SetAccent{149A98}{EAF7F6}
\CompactHero{From innovation capability to outcomes}{What the three pillars are designed to enable}{The objective is not idea generation for its own sake. Students learn to convert differentiated, substantive and useful ideas into stronger research, invention and industry-facing evidence.}

\begin{BroCard}
\SectionTitle{Outcome 01 - Impactful research}
\begin{minipage}[t]{.31\linewidth}\SmallCaps{Uniqueness}\par\TightText{A consequential question, clearer research gap and a contribution that is meaningfully different.}\end{minipage}\hfill
\begin{minipage}[t]{.31\linewidth}\SmallCaps{Hard to do}\par\TightText{Technical depth, rigorous methods, experiments, difficult trade-offs and credible evidence.}\end{minipage}\hfill
\begin{minipage}[t]{.31\linewidth}\SmallCaps{Value adding}\par\TightText{A problem that matters to a field, stakeholder, industry, society or future system.}\end{minipage}
\vspace{2.3mm}
\begin{SoftCard}\SmallCaps{Student-owned evidence}\par\vspace{.8mm}\TightText{Problem landscape \textbullet\ novelty map \textbullet\ research question \textbullet\ experiment plan \textbullet\ contribution statement \textbullet\ publication pathway.}\end{SoftCard}
\end{BroCard}

\CardGap
\begin{BroCard}
\SectionTitle{Outcome 02 - Patents and defensible IP}
\begin{minipage}[t]{.31\linewidth}\SmallCaps{Uniqueness}\par\TightText{Prior-art awareness and a clearly differentiated mechanism, method or system element.}\end{minipage}\hfill
\begin{minipage}[t]{.31\linewidth}\SmallCaps{Hard to do}\par\TightText{Non-obvious technical substance, constraints and implementation depth that make the idea defensible.}\end{minipage}\hfill
\begin{minipage}[t]{.31\linewidth}\SmallCaps{Value adding}\par\TightText{Utility, adoption potential and a reason the invention should exist beyond novelty alone.}\end{minipage}
\vspace{2.3mm}
\begin{SoftCard}\SmallCaps{Student-owned evidence}\par\vspace{.8mm}\TightText{Invention canvas \textbullet\ prior-art search terms \textbullet\ novelty hypothesis \textbullet\ non-obviousness rationale \textbullet\ prototype or evidence plan.}\end{SoftCard}
\end{BroCard}

\CardGap
\begin{BroCard}
\SectionTitle{Outcome 03 - Industry-ready innovators}
\begin{minipage}[t]{.31\linewidth}\SmallCaps{Uniqueness}\par\TightText{Spot opportunities others miss and frame a differentiated response rather than repeating obvious solutions.}\end{minipage}\hfill
\begin{minipage}[t]{.31\linewidth}\SmallCaps{Hard to do}\par\TightText{Move from idea to execution: handle constraints, collaborate, experiment, learn and improve.}\end{minipage}\hfill
\begin{minipage}[t]{.31\linewidth}\SmallCaps{Value adding}\par\TightText{Connect technical work to business, user and societal outcomes and explain why it matters.}\end{minipage}
\vspace{2.3mm}
\begin{SoftCard}\SmallCaps{Student-owned evidence}\par\vspace{.8mm}\TightText{Problem brief \textbullet\ prototype or concept \textbullet\ decision rationale \textbullet\ outcome story \textbullet\ pitch \textbullet\ 30/60-day action plan.}\end{SoftCard}
\end{BroCard}

\CardGap
\begin{DarkCard}
{\color{BroGold}\fontsize{8.2}{9.8}\selectfont\bfseries WHAT ACADEMIC LEADERS CAN REVIEW AFTER A COHORT\par}\vspace{1.4mm}
{\color{white}\fontsize{8.9}{10.9}\selectfont Problem portfolio \textbullet\ research directions \textbullet\ possible IP candidates \textbullet\ experiment plans \textbullet\ industry-facing pitches \textbullet\ participant capability evidence \textbullet\ recommended next steps.}
\end{DarkCard}
\vspace{2mm}
{\fontsize{7.5}{9}\selectfont\color{BroMuted}\textit{The program builds capability and stronger pathways. Publication, patent and hiring outcomes are not guaranteed and depend on the quality of subsequent work, evidence, review and execution.}}
\ProgramFooter{FROM BETTER QUESTIONS TO STRONGER RESEARCH, IP AND INDUSTRY READINESS}
\newpage
